\documentclass[10pt,a4paper]{amsart}
\usepackage[margin=19mm]{geometry}
\usepackage{amsmath,amssymb,mathtools}
\usepackage[scr=rsfs,cal=euler]{mathalfa}
\usepackage{xfrac}
\usepackage[unicode,hidelinks,bookmarks=false]{hyperref}
\newcommand{\ie}{i.e.\ }
\newcommand{\cf}{cf.\ }
\newcommand{\resp}{respectively\ }
\newcommand{\Subsection}{\subsection}
\providecommand{\nobreakdash}{\nobreak\hbox{-}}
\setlength{\emergencystretch}{2em}
\multlinegap0pt
\usepackage{bm}
\usepackage{bbm}
\usepackage{dsfont}
\usepackage{xspace}
\usepackage{cancel}
\usepackage{mathtools}
\usepackage{amsthm}
\makeatletter
\@ifundefined{lemma}{\newtheorem{lemma}{Lemma}}{}
\makeatother
\newcommand{\Lip}{\mathrm {Lip}}
\newcommand{\cF}{\mathcal {F}}
\newcommand{\dX}{\mathbb{X}}
\newcommand{\diam}{\mathrm{diam}}
\newcommand{\dist}{\mathrm{dist}}
\title[Modern aspects of Markov chains: entropy, curvature and the ]{Modern aspects of Markov chains: entropy, curvature and the cutoff phenomenon}
\author{Justin Salez}
\date{Source: arXiv:2508.21055v1 (CC BY 4.0)}
\begin{document}
\maketitle

\noindent\emph{Source and licence.} Modern aspects of Markov chains: entropy, curvature and the cutoff phenomenon. Version arXiv:2508.21055v1 (CC BY 4.0), distributed under the Creative Commons Attribution 4.0 International licence, as recorded in the arXiv licence metadata for this exact version and shown on the abstract page https://arxiv.org/abs/2508.21055v1. Full rights notice: licenses/arxiv-2508.21055-CC-BY-4.0.txt.\par\smallskip

\noindent\textit{Editorial scope.} Counted unit: the Lemma whose source label is \texttt{lm:lip} in the article (source0.tex lines 2220--2225, source label \texttt{lm:lip}), delivered together with the article's own complete original proof of it (source0.tex lines 2226--2265, one \texttt{proof} environment of 40 lines). No mathematics is written, repaired or re-derived: statement and proof are the article's own text line for line. Required context retained uncounted: none. Every definition, display and label of those lines is retained unchanged. Omitted from this package and not redistributed: 1--2219, 2266--2481. Nothing inside a retained excerpt is omitted or altered: every retained body is delivered exactly as the article's lines are written, so no omission receipt of any kind accompanies this package. Citations: the retained excerpts carry no citation command at all, so no bibliography entry of the article is transcribed and no reference is redistributed. Reference adaptation: none was needed and none was made; every reference inside the delivered unit resolves inside it, so no unresolved reference or citation is produced. Printed slips kept literal: no printed word, symbol, hypothesis or number of the counted statement or of its proof was altered. Layout and class: the article's own class is \texttt{article} with packages including amsmath, amsthm, amsfonts, amssymb, graphicx, hyperref, enumerate, psfrag, tikz, pgfplots, mathrsfs, xcolor, fullpage, inputenc; the wrapper is the lane's pinned \texttt{amsart} editorial wrapper at 10pt on A4, which restates the theorem-like environments the retained text needs and adds the article's own macro definitions that the retained text needs (\texttt{\textbackslash Lip}, \texttt{\textbackslash cF}, \texttt{\textbackslash dX}, \texttt{\textbackslash diam}, \texttt{\textbackslash dist}), with extra wrapper packages (bm, bbm, dsfont, xspace, cancel, mathtools). The delivered numbering is the unit's own. Assets: no retained span contains a figure environment, a graphics-inclusion command, a PDF-page inclusion, a table or a TikZ picture; no image file is copied into this package, the package declares no asset dependency and no image file is redistributed with it. Licence: version arXiv:2508.21055v1 of this article is distributed under the Creative Commons Attribution 4.0 International licence, as recorded in the arXiv licence metadata for this exact version and shown on the abstract page https://arxiv.org/abs/2508.21055v1. A full rights notice is delivered at licenses/arxiv-2508.21055-CC-BY-4.0.txt. Characters: every retained excerpt is pure ASCII, so the delivered engine, XeLaTeX with its default Latin Modern text font, reports no missing character.
\begin{lemma}[Universal roughness estimate]\label{lm:lip}For any non-zero $f\in \cF^+$ and any $t> 0$,
\begin{eqnarray*}
\Lip \log P_t f & \le & 3\log d+2\log_+\left(\frac{\diam(\dX)}{t}\right).
\end{eqnarray*}
Moreover, the same conclusion applies to $P_t^\star$ instead of $P_t$, with $d$ replaced by $d^2$.
\end{lemma}

\begin{proof}
Fix $t>0$ and set $p_k:=e^{-t}t^k/k!$, so that we can write
\begin{eqnarray*}
P_t & = & \sum_{k=0}^\infty p_k T^{k} \ = \ \frac{1}{t}\sum_{k=1}^\infty {kp_{k}}T^{k-1}, 
\end{eqnarray*}
thanks to the Poisson identity $tp_k=(k+1)p_{k+1}$ and a change of indices. In particular, 
\begin{eqnarray*}
\frac{TP_t f}{P_tf} & = & \frac{1}{t}\frac{\sum_{k=0}^\infty k p_{k} T^{k}f}{\sum_{k=0}^\infty p_{k} T^{k}f} \
 \le \ \frac{1}{t}\log\left(\frac{\sum_{k=0}^\infty p_{k}e^k T^{k}f}{\sum_{k=0}^\infty p_{k} T^{k}f}\right),
 \end{eqnarray*}
by Jensen's inequality. Let us now estimate the ratio inside the logarithm. For the numerator, we simply use the crude bound $T^kf\le \|T^k f\|_\infty\le \|f\|_\infty$ to obtain
\begin{eqnarray*}
\sum_{k=0}^\infty p_{k}e^k T^{k}f & \le & \|f\|_\infty e^{(e-1)t},
\end{eqnarray*}
while for the denominator, we write
\begin{eqnarray*}
\sum_{k=0}^\infty p_{k} T^{k}f & \ge & \|f\|_\infty \min_{(x,y)\in\dX^2}\left\{\sum_{k=0}^\infty p_{k} T^{k}(x,y)\right\},
\end{eqnarray*}
and  estimate the minimum on the right-hand side in the following crude way: given $(x,y)\in\dX^2$, we consider a path $(x_0,\ldots,x_\ell)$ from $x_0=x$ to $x_\ell=y$ of length  $\ell=\dist(x,y)$. Using the definition of $d$  and the classical bound  $\ell!\le \ell^\ell$ (with $0^0=1$), we have
\begin{eqnarray*}
\sum_{k=0}^\infty p_kT^k(x,y) & \ge & p_\ell T(x_0,x_1)\cdots T(x_{\ell-1},x_\ell)\ \ge \ e^{-t}\left(\frac{t}{d\ell}\right)^\ell,
 \end{eqnarray*} 
and since $u\mapsto u\log u$ is convex, the right-hand side is minimized when $\ell=0$ or $\ell=\diam(\dX)$. In view of the last four displays, we conclude that
\begin{eqnarray*}
\left\|\frac{TP_t f}{P_tf} \right\|_\infty & \le & e+\frac{\diam(\dX)}{t}\log_+\left(\frac{d\diam(\dX)}{t}\right).
 \end{eqnarray*}
On the other hand,  we  have by definition
\begin{eqnarray*}
\left\|\frac{TP_t f}{P_tf} \right\|_\infty & = & \max_{x\in\dX}\left\{\sum_{y\in\dX}\frac{T(x,y)P_tf(y)}{P_tf(x)}\right\} \ \ge \  \max_{(x,y)\in E}\left\{\frac{T(x,y)P_tf(y)}{P_tf(x)}\right\} \ \ge \ \frac{e^{\Lip\log P_tf}}{d}.
\end{eqnarray*}
Combining those two bounds yields 
\begin{eqnarray*}
\Lip\log P_tf & \le & \log\left[de+\frac{d\diam(\dX)}{t}\log_+\left(\frac{d\diam(\dX)}{t}\right)\right],
\end{eqnarray*}
which easily leads to the main claim. Finally,  the identity 
\begin{eqnarray*}
\forall (x,y)\in\dX^2,\qquad T^\star(x,y)T^\star(y,x) & = & T(x,y)T(y,x),
\end{eqnarray*}
shows that the non-zero entries of $T^\star$ are at least $d^{-2}$, hence the second claim. 
\end{proof}

\end{document}
